\documentclass[11pt]{article}

\usepackage[margin=1in]{geometry}
\usepackage{amsmath,amssymb,amsthm}
\usepackage{graphicx}
\usepackage[hidelinks]{hyperref}

\newtheorem{theorem}{Theorem}
\newtheorem{corollary}{Corollary}
\newtheorem{remark}{Remark}

\newcommand{\HH}{\mathcal H}
\newcommand{\BB}{\mathcal B}
\newcommand{\NN}{\mathcal N}
\newcommand{\RR}{\mathcal R}
\newcommand{\opRe}{\operatorname{Re}}

\title{A nilpotent pure-part result for the Fong--Tsui problem}
\author{Run \texttt{fa\_banach\_001}, target arXiv:1511.09055}
\date{2026-07-03}

\begin{document}
\maketitle

\paragraph{Status.}
Partial result, likely valid. The result answers the finite nilpotent
obstruction which appears in Mbekhta--Suciu's analysis of the Fong--Tsui
conjecture for contractive \(m\)-quasi-isometries. It does not settle the
general Fong--Tsui conjecture for arbitrary pure contractions.

\section*{Source question}

The source is Mostafa Mbekhta and Laurian Suciu,
\emph{Partial isometries and the conjecture of C. K. Fong and S. K. Tsui},
arXiv:1511.09055 \cite{MbekhtaSuciu}. They study bounded operators \(T\) on a
complex Hilbert space satisfying
\begin{equation}\label{eq:FT}
  |T|\leq |\opRe T|,
  \qquad |T|=(T^*T)^{1/2}, \quad \opRe T=\frac{T+T^*}{2}.
\end{equation}
After proving the conjecture for contractive \(2\)-quasi-isometries, they ask
on page 11 whether the hypothesis \(Q^2=0\) in Proposition 3.3 can be replaced
by \(Q^m=0\) for some \(m\geq 3\), in order to preserve the conclusion \(Q=0\)
under condition \eqref{eq:FT}.

\begin{center}
\includegraphics[width=\textwidth]{figures/open_problem_crop.png}
\end{center}

\section*{Proof intuition}

The finite nilpotent case has a numerical-radius obstruction which is not used
in the source proof. If \(N^m=0\), then the Haagerup--de la Harpe theorem gives
a strict upper bound
\[
  w(N)\leq \|N\|\cos\frac{\pi}{m+1}<\|N\|
\]
for nonzero \(N\), after normalizing. On the other hand, the Fong--Tsui
inequality \(|N|\leq |\opRe N|\) forces
\[
  \|N\|=\||N|\|\leq \||\opRe N|\|=\|\opRe N\|\leq w(N).
\]
Thus a nonzero nilpotent would have numerical radius simultaneously at least
and strictly less than its norm. The only possibility is \(N=0\).

\section*{Main result}

\begin{theorem}[Finite nilpotent blocks vanish]\label{thm:nilpotent}
Let \(N\in\BB(\HH)\) be a bounded operator on a complex Hilbert space. If
\(N^m=0\) for some integer \(m\geq 2\) and
\[
  |N|\leq |\opRe N|,
\]
then \(N=0\).
\end{theorem}

\begin{proof}
Assume \(N\neq 0\). Since \(0\leq |N|\leq |\opRe N|\), monotonicity of the norm
on positive operators gives
\[
  \|N\|=\||N|\|\leq \||\opRe N|\|=\|\opRe N\|.
\]
For a self-adjoint operator \(A\), \(\|A\|=\sup_{\|x\|=1}|\langle Ax,x\rangle|\).
Applying this to \(A=\opRe N\),
\[
  \|\opRe N\|
  = \sup_{\|x\|=1} |\operatorname{Re}\langle Nx,x\rangle|
  \leq \sup_{\|x\|=1} |\langle Nx,x\rangle|
  = w(N),
\]
where \(w(N)\) is the numerical radius.

The Haagerup--de la Harpe numerical-radius inequality for nilpotents
\cite{HaagerupDeLaHarpe} says that if \(N^m=0\), \(m\geq 2\), then
\[
  w(N)\leq \|N\|\cos\frac{\pi}{m+1}.
\]
Combining the inequalities yields
\[
  \|N\|\leq w(N)\leq \|N\|\cos\frac{\pi}{m+1}.
\]
Since \(0<\cos(\pi/(m+1))<1\), this contradicts \(\|N\|>0\). Hence \(N=0\).
\end{proof}

\begin{corollary}[Replacement in Proposition 3.3]\label{cor:prop33}
In Proposition 3.3 of \cite{MbekhtaSuciu}, the alternative hypothesis
\(Q^2=0\) may be replaced by \(Q^m=0\) for any finite \(m\geq 2\). More
precisely, let
\[
  T=\begin{pmatrix} S&R\\0&Q\end{pmatrix}
  \quad\text{on}\quad \HH=\mathcal G\oplus\mathcal G^\perp,
\]
where \(\|T\|^{-1}S\) is an isometry on \(\mathcal G\), and assume
\(Q^m=0\) for some finite \(m\geq 2\). If \(T\) satisfies \eqref{eq:FT}, then
\(T\) is self-adjoint.
\end{corollary}

\begin{proof}
If \(T=0\), there is nothing to prove. Otherwise scale as in the proof of
\cite[Proposition 3.3]{MbekhtaSuciu}, writing
\(T_0=\|T\|^{-1}T\), \(S_0=\|T\|^{-1}S\), \(R_0=\|T\|^{-1}R\), and
\(Q_0=\|T\|^{-1}Q\). The initial part of that proof uses Theorem 3.1 of the
source paper and does not use \(Q^2=0\). It gives
\[
  R_0R_0^*=P_{\NN(S_0^*)},\qquad Q_0R_0^*=0,
\]
and, in the nilpotent alternative, that
\(\NN(I-Q_0Q_0^*)\) is invariant under \(Q_0^*\).

If \(b\in\NN(I-Q_0Q_0^*)\), then \(b=Q_0Q_0^*b\). Since this nullspace is
\(Q_0^*\)-invariant, an induction gives
\[
  b=Q_0^kQ_0^{*k}b \qquad (1\leq k\leq m).
\]
Taking \(k=m\) and using \(Q_0^m=0\), we obtain \(b=0\). Hence
\(\NN(I-Q_0Q_0^*)=\{0\}\). The remainder of the reduction in the source proof
then yields \(R_0=0\), so \(T_0=S_0\oplus Q_0\), and both summands satisfy
\eqref{eq:FT}.

Theorem \ref{thm:nilpotent} applied to \(Q_0\) gives \(Q_0=0\). The same source
argument gives that \(S_0\) is a symmetry. Thus
\(T=\|T\|(S_0\oplus 0)\) is self-adjoint.
\end{proof}

\begin{remark}[Blockwise inheritance]
The last step uses the elementary fact that if \(T=S\oplus Q\) is a reducing
direct sum and \(T\) satisfies \eqref{eq:FT}, then \(Q\) also satisfies
\eqref{eq:FT}. Indeed,
\[
  |S\oplus Q|=|S|\oplus |Q|,
  \qquad
  |\opRe(S\oplus Q)|=|\opRe S|\oplus |\opRe Q|.
\]
Restricting the positive operator inequality to the \(Q\)-summand gives
\(|Q|\leq |\opRe Q|\).
\end{remark}

\begin{corollary}[Contractive finite \(m\)-quasi-isometries]
Let \(T\in\BB(\HH)\) be a contractive \(m\)-quasi-isometry for a finite integer
\(m\geq 1\). If \(T\) satisfies \eqref{eq:FT}, then \(T\) is self-adjoint.
\end{corollary}

\begin{proof}
For \(m=1\) and \(m=2\), this is already contained in the source paper: see the
quasi-isometry case and Corollary 3.4 of \cite{MbekhtaSuciu}. Let \(m\geq 3\).
Corollary 3.5 of \cite{MbekhtaSuciu} shows that
\[
  \HH=\overline{\RR(T^m)}\oplus \NN(T^{*m})
\]
reduces \(T\) to the block diagonal form
\[
  T=S\oplus Q,
\]
where \(S\) is a symmetry on \(\overline{\RR(T^m)}\) and \(Q^m=0\). Since
\eqref{eq:FT} holds for \(T\), the preceding blockwise inheritance gives
\(|Q|\leq |\opRe Q|\). Theorem \ref{thm:nilpotent} gives \(Q=0\). Hence
\(T=S\oplus 0\), so \(T=T^*\).
\end{proof}

\section*{Verification notes and limitations}

The proof is analytic and uses no finite search. The two main points for review
are:
\begin{enumerate}
\item the Haagerup--de la Harpe bound applies to arbitrary bounded
Hilbert-space nilpotents after scaling by \(\|N\|\);
\item in a reducing direct sum \(T=S\oplus Q\), the operator inequality
\(|T|\leq |\opRe T|\) restricts to each summand.
\end{enumerate}

This packet does not prove the full Fong--Tsui conjecture and does not address
general pure contractions whose pure part is not finite nilpotent. Corollary
\ref{cor:prop33} closes the finite nilpotent obstruction left open on page 11
of the source paper and upgrades the source's contractive \(m\)-quasi-isometry conclusion from
``\(T\) is self-adjoint if and only if \(T=S\oplus 0\)'' to self-adjointness for
every finite \(m\).

\section*{Novelty check}

Local run indexes were searched for arXiv:1511.09055, the title terms,
``Fong--Tsui'', \(Q^m=0\), nilpotent pure part, and Haagerup numerical radius;
no existing packet or attempt for this target was found. Web searches for
``Fong-Tsui conjecture nilpotent'', ``Fong-Tsui \(Q^m=0\)'',
``\(|T|\leq |\operatorname{Re}T|\) nilpotent'', and the Haagerup--de la Harpe
nilpotent numerical-radius theorem found the source paper, Mortad's earlier
contribution, and numerical-radius references, but no later explicit answer to
the page-11 nilpotent question.

\begin{thebibliography}{9}

\bibitem{MbekhtaSuciu}
M. Mbekhta and L. Suciu,
\emph{Partial isometries and the conjecture of C. K. Fong and S. K. Tsui},
arXiv:1511.09055.

\bibitem{HaagerupDeLaHarpe}
U. Haagerup and P. de la Harpe,
\emph{The numerical radius of a nilpotent operator on a Hilbert space},
Proceedings of the American Mathematical Society \textbf{115} (1992),
371--379.

\bibitem{Mortad}
M. H. Mortad,
\emph{A contribution to the Fong--Tsui conjecture related to self-adjoint
operators}, arXiv:1208.4346.

\end{thebibliography}

\end{document}
